\section{第四天}
\begin{enumerate}
  \item 核心要义之\bs{建设法治体系}（全面依法治国的总抓手）。
        \begin{enumerate}
          \item 全推国，\bs{总体目标}是\bs{建法系，建法国}。
          \item 这就是，在中国共产\bs{党领导}下，\bs{坚持中特制}，\bs{贯彻中法论}，
          \item 法治体系。\zx{形成五个体系}，\bs{完备的法律规范体系、高效的法治实施体系、严密的法治监督体系、有力的法治保障体系、完善的党内法规体系}，
          \item 法治国家。
                \begin{enumerate}
                  \item 坚持\zx{三依共同推进}，\bs{依法治国、依法执政、依法行政}。
                  \item 坚持\zx{三治一体建设}，\bs{法治国家、法治政府、法治社会}。
                  \item 实现\zx{立执司守带修饰}，\bs{科学立法、严格执法、公正司法、全民守法}。
                  \item 促进\zx{一个现代化}。\bs{国家治理体系和治理能力现代化}。
                \end{enumerate}
          \item 【说明】全面依法治国的总目标已经包含了对法治体系的说明，因为总目标有可能单独考查，因此可以合并背诵。
                \begin{enumerate}
                  \item \zx{法治体系为}五个体系。
                  \item \zx{法治国家为}坚持三依共同推进，坚持三治一体建设，
                        实现立执司守带修饰，促进一个现代化（三依三治，四项工作，一个现代化）。
                \end{enumerate}
        \end{enumerate}
  \item 核心要义之坚持\bs{在法治轨道上推进国家治理体系和治理能力现代化}。
        \begin{enumerate}
          \item 概念和意义（深化改革的总目标）。
                \begin{enumerate}
                  \item 治理体系。
                        \begin{enumerate}
                          \item 国家治理体系是在\zx{党领导}下\zx{管理国家}的\bs{制度体系}，
                          \item 包括\zx{五位一体}等各领域的\zx{体制机制、法律法规安排}，\bs{经济、政治、文化、社会、生态文明和党的建设}。
                          \item 是一整套\zx{紧密相连、相互协调}的\bs{国家制度}。
                        \end{enumerate}
                  \item 治理能力。
                        \begin{enumerate}
                          \item 国家治理能力是\bs{运用国家制度}\zx{管理社会各方面事务}的能力，
                          \item 是\zx{改革发展稳定、内政外交国防、治党治国治军}等各个方面\bs{国家制度}\zx{执行能力的集中体现}。
                        \end{enumerate}
                  \item 法治轨道。
                        \begin{enumerate}
                          \item 一现化，是国家\zx{治理现代化的具体表现}，
                          \item 其\zx{核心内涵}主要是实现国家\zx{治理的制度化、程序化、法治化}。
                          \item 因此，\zx{法治}是国家治理体系和治理能力的\zx{重要依托}，
                          \item 总结。坚持\zx{在法轨上}推进一现化，\bs{是实现良法善治的必由之路}。
                        \end{enumerate}
                  \item 价值意义。万能价值。
                \end{enumerate}
          \item 如何。
                \begin{enumerate}
                  \item 更好发挥法治\zx{固根本、稳预期、利长远}的保障作用。
                  \item 坚持依法\zx{治军}，从严治军。
                  \item 坚持依法\zx{治网}。
                  \item 坚持依法保障\zx{“一国两制”实践，推进祖国统一}。
                \end{enumerate}
        \end{enumerate}
\end{enumerate}